%%%%%%%%%%  scan idx 48  |  folio 35  %%%%%%%%%%
au
sous-anneau de $K$ engendré par les $\lambda^iN$ pour $1\leq i\leq q$ et par l'inverse de $\lambda^qN$. D'autre part, on vérifie facilement dans cet anneau que l'on a $\ell^p(N)\lambda_{-1}(N)=\lambda^p(\lambda_{-1}(N))$. Comme on est dans un grand corps $K$ et que $\lambda_{-1}(N)\neq0$, on en conclut $\ell^p(N)=\lambda^p(N,1)$~; les deux membres sont dans l'intersection de $\underline{K}(G)$ et de l'anneau des séries formelles en les $\lambda^iN$, donc sont des polynômes en les $\lambda^iN$.

Pour prouver (1.33) dans le cas général, on vérifie que le deuxième membre satisfait dans $\underline{K}(G)$ à la formule analogue à (1.32), c'est-à-dire que son produit par $\lambda_{-1}(N)$ est $\lambda^p(F\lambda_{-1}(N))$. De même, le produit de $\lambda^p(N,F)\in \underline{K}(G)$ par $\lambda_{-1}(N)$ est $\lambda^p(F\lambda_{-1}(N))$, comme il résulte du fait que $\underline{K}(G)$ est \emph{spécial}. Comme $\underline{K}(G)$ est intègre (cf \S 2) et que $\lambda_{-1}(N)\neq0$ (même référence), on a bien démontré la formule (1.33). cqfd.

\textbf{Remarques :} a) Lorsque la dimension $q'$ de $F$ est $\geq p$, alors (1.33) caractérise le polynôme $\lambda^p(N,x)$ en les $\lambda^iN$ ($1\leq i\leq q$) et les $\lambda^jx$ ($1\leq j\leq p$).

b) Il est plausible que (1.33) est encore vérifiée si le corps $k$ n'est pas algébriquement clos\footnote{Cela se prouve en effet par la même méthode.}. Il faudrait résoudre cette question si l'on voulait prouver le théorème de Riemann--Roch sur un corps de caractéristique $0$ non algébriquement clos par la méthode qu'on va exposer. Le passage d'un $\mathfrak S_p$-module à sa partie alternée n'a d'ailleurs de sens raisonnable que si le corps de base est de caractéristique $0$.

Soit maintenant $K$ un $\lambda$-anneau \emph{spécial} quelconque et soit $N$ un
élé%
